% Storyline 1 (final), paragraph 1: classical result, the two steps of reputation formation, the central problem.
Cooperation among unrelated individuals is sustained in part by reputation\cite{milinski2016,nowak2005}. By helping others, individuals improve their standing and become more likely to receive help from third parties, a mechanism known as indirect reciprocity\cite{nowak1998,leimar2001,wedekind2000,milinski2002}. Reputations are built in two steps. Observers judge each action against a social norm, and they revise their view of the actor in light of that judgement\cite{ohtsuki2004,santos2018}. Theory has identified norms that make reputation-based cooperation evolutionarily stable\cite{panchanathan2003,ohtsuki2004}, most notably the leading-eight norms, which sustain cooperation when reputations are public, even in the presence of occasional errors\cite{ohtsuki2004,ohtsuki2006}. These results rest on agreement about who is good and who is bad. When reputations are private and assessments are noisy, individuals hold different views of the same person\cite{takahashi2006,uchida2010}, and most leading-eight norms no longer maintain stable cooperation\cite{hilbe2018,okada2018,perret2021,fujimoto2023}. Reputations must then remain both accurate and shared, distinguishing those who cooperate from those who do not while different observers largely agree about who is who\cite{fujimoto2022,murase2024synchronization}.
